\section{Цель работы}

Использование численных методов для решения задачи Коши обыкновенных дифференциальных уравнений.

\section{Описание алгоритма решения задачи}

Задача Коши для обыкновенного дифференциального уравнения первого порядка: требуется найти функцию $Y = Y(x)$, удовлетворяющую уравнению $Y' = f(x, Y)$ и начальному условию $Y(x_0) = Y_0$. Численное решение задачи Коши состоит в приближённом вычислении значений $y_0, y_1, \dots, y_n$ искомой функции в узлах сетки $x_0, x_1, \ldots, x_n$. Метод конечных разностей заключается в замене области непрерывного изменения аргумента дискретной сеткой узлов и аппроксимации дифференциального уравнения разностным.

В данной работе (вариант 14) реализованы три метода:

\begin{itemize}
\item \textbf{Модифицированный метод Эйлера} — одношаговый явный метод. На каждом шаге сначала вычисляется грубое приближение по формуле Эйлера, затем оно уточняется путём усреднения значений правой части в текущем и следующем узлах. Порядок точности $O(h^2)$.

\item \textbf{Метод Рунге-Кутта четвёртого порядка} — одношаговый явный метод. На каждом шаге производится четырёхкратное вычисление правой части $f(x, y)$ в промежуточных точках, после чего вычисляется итоговое приращение как взвешенная сумма полученных значений. Порядок точности $O(h^4)$.

\item \textbf{Метод Милна} — многошаговый метод прогноза и коррекции (предиктор-корректор). На этапе предиктора по явной формуле вычисляется начальное приближение, на этапе корректора оно уточняется по неявной формуле. Для запуска метода первые три точки вычисляются методом Рунге-Кутта. Порядок точности $O(h^4)$.
\end{itemize}

Для оценки погрешности одношаговых методов используется правило Рунге. Для многошагового метода — сравнение с точным решением.

\section{Блок-схемы алгоритмов}

\subsection{Блок-схема модифицированного метода Эйлера}

Здесь размещается блок-схема модифицированного метода Эйлера.

\subsection{Блок-схема метода Рунге-Кутта четвёртого порядка}

Здесь размещается блок-схема метода Рунге-Кутта четвёртого порядка.

\subsection{Блок-схема метода Милна (предиктор-корректор)}

Здесь размещается блок-схема метода Милна (метода предиктор-корректор).

\section{Рабочие формулы метода}

\subsection{Модифицированный метод Эйлера (метод Эйлера с пересчётом)}

Модифицированный метод Эйлера является одношаговым явным методом. На каждом шаге сначала вычисляется грубое приближение по формуле Эйлера, затем оно уточняется. Расчётные формулы:

\[
\tilde{y}_{i+1} = y_i + h f(x_i, y_i)
\]

\[
\mathbf{y}_{i+1} = \mathbf{y}_i + \frac{h}{2}\left[ \mathbf{f}\left(\mathbf{x}_i, \mathbf{y}_i\right) + \mathbf{f}\left(\mathbf{x}_{i+1}, \mathbf{y}_i + \mathbf{h f}\left(\mathbf{x}_i, \mathbf{y}_i\right)\right) \right], \quad i = 0, 1 \dots
\]

Начальное значение $\mathbf{y}_0 = \mathbf{Y}_0$ задано начальным условием.

Порядок точности метода: $\delta_n = O(h^2)$.

\subsection{Метод Рунге-Кутта четвёртого порядка}

Метод Рунге-Кутта четвёртого порядка является одношаговым явным методом, требующим на каждом шаге четырёхкратного вычисления правой части $f(x, y)$. Расчётные формулы:

\[
\mathbf{y}_{i+1} = \mathbf{y}_i + \frac{1}{6}(\mathbf{k}_1 + 2\mathbf{k}_2 + 2\mathbf{k}_3 + \mathbf{k}_4),
\]

\[
\mathbf{k}_1 = \mathbf{h} \cdot \mathbf{f}(\mathbf{x}_i, \mathbf{y}_i)
\]

\[
\mathbf{k}_2 = \mathbf{h} \cdot \mathbf{f}\left(x_i + \frac{h}{2}, y_i + \frac{k_1}{2}\right)
\]

\[
\mathbf{k}_3 = \mathbf{h} \cdot \mathbf{f}\left(x_i + \frac{\mathbf{h}}{2}, y_i + \frac{\mathbf{k}_2}{2}\right)
\]

\[
\mathbf{k}_4 = \mathbf{h} \cdot \mathbf{f}(x_i + \mathbf{h}, y_i + \mathbf{k}_3)
\]

Порядок точности метода: $\delta_n = O(h^4)$.

\subsection{Метод Милна (предиктор-корректор)}

Метод Милна относится к многошаговым методам и представляет один из методов прогноза и коррекции. Для получения формул Милна используется первая интерполяционная формула Ньютона с разностями до третьего порядка. Решение в следующей точке находится в два этапа.

а) Этап прогноза:

\[
y_i^{\text{прогн}} = y_{i-4} + \frac{4h}{3}(2f_{i-3} - f_{i-2} + 2f_{i-1})
\]

б) Этап коррекции:

\[
\begin{array}{l}
y_i^{\text{корр}} = y_{i-2} + \dfrac{h}{3}(f_{i-2} + 4f_{i-1} + f_i^{\text{прогн}}) \\[6pt]
f_i^{\text{прогн}} = f(x_i, y_i^{\text{прогн}})
\end{array}
\]

Для начала счёта требуется задать решения в трёх первых точках, которые получаются одношаговым методом (в программе — методом Рунге-Кутта).

Порядок точности метода: $\delta_n = O(h^4)$.

\subsection{Оценка погрешности}

Для одношаговых методов (модифицированный метод Эйлера и метод Рунге-Кутта) оценка погрешности производится по правилу Рунге:

\[
R = \frac{\left| y^h - y^{h/2} \right|}{2^p - 1},
\]

где $y^h$ — решение задачи Коши с шагом $h$ в точке $x + h$, $y^{h/2}$ — решение задачи Коши с шагом $h/2$ в точке $x + h$, $p$ — порядок точности метода ($p = 2$ для модифицированного метода Эйлера, $p = 4$ для метода Рунге-Кутта).

Для многошагового метода Милна оценка погрешности производится путём сравнения с точным решением:

\[
\varepsilon = \max_{0 \leq i \leq n} \left| y_{i \text{ точн}} - y_i \right|,
\]

где $y_{i \text{ точн}}$, $y_i$ — значения точного и приближенного решений в узлах сетки $x_i$, $i = 1, \dots, n$.
